\section{Invariants}

\begin{frame}
\centering
\Huge\textbf{Invariants}
\end{frame}

% --------------------------------------------------
% Question 1
% --------------------------------------------------

\begin{frame}{Question 1}
A program starts with an integer counter $x=0$.

At every step, the program either increases $x$ by $2$ or decreases $x$ by $2$.

\textbf{Question:} Identify an invariant of the counter and explain why it remains true after every operation.
\end{frame}

\begin{frame}{Answer 1}
Initially,

$$
x=0,
$$

which is even.

Suppose $x$ is even before an operation.

The program changes $x$ to either

$$
x+2
$$

or

$$
x-2.
$$

Both $2$ and $-2$ are even.
\end{frame}

\begin{frame}{Answer 1}
Therefore,

$$
\text{even}+\text{even}=\text{even}.
$$

Hence, after either operation, $x$ remains even.

Thus, the invariant is

$$
\boxed{x\text{ is always even}.}
$$

\end{frame}

% --------------------------------------------------
% Question 2
% --------------------------------------------------

\begin{frame}{Question 2}
A token is placed on position $0$ of a one-dimensional grid.

At each step, the token moves either $2$ positions to the left or $2$ positions to the right.

\textbf{Question:} Can the token ever reach position $5$? Use an invariant to justify your answer.
\end{frame}

\begin{frame}{Answer 2}
The token starts at

$$
x=0,
$$

which is even.

Each move changes the position by either

$$
+2
\quad\text{or}\quad
-2.
$$

Therefore, the parity of the position never changes.

The invariant is

$$
\boxed{x\text{ is always even}.}
$$

\end{frame}

\begin{frame}{Answer 2}
Since $5$ is odd,

$$
5\not\equiv0\pmod 2.
$$

But every reachable position is even.

Therefore, position $5$ cannot be reached.

$$
\boxed{\text{The token can never reach position }5.}
$$

\end{frame}

% --------------------------------------------------
% Question 3
% --------------------------------------------------

\begin{frame}{Question 3}
A distributed system has two counters $A$ and $B$.

Initially,

$$
A+B=20.
$$

At every operation, the system increases $A$ by $1$ and decreases $B$ by $1$, or decreases $A$ by $1$ and increases $B$ by $1$.

\textbf{Question:} Determine an invariant of the system.
\end{frame}

\begin{frame}{Answer 3}
Initially,

$$
A+B=20.
$$

Consider the first type of operation:

$$
A\rightarrow A+1,
\qquad
B\rightarrow B-1.
$$

The new sum is

$$
(A+1)+(B-1)=A+B.
$$

\end{frame}

\begin{frame}{Answer 3}
For the second operation,

$$
A\rightarrow A-1,
\qquad
B\rightarrow B+1.
$$

The new sum is

$$
(A-1)+(B+1)=A+B.
$$

Therefore, neither operation changes the sum.

Hence,

$$
\boxed{A+B=20}
$$

is an invariant of the system.
\end{frame}

% --------------------------------------------------
% Question 4
% --------------------------------------------------

\begin{frame}{Question 4}
A program contains two integer variables $x$ and $y$.

Initially,

$$
x+y=10.
$$

The program repeatedly replaces $(x,y)$ by either

$$
(x+1,y-1)
$$

or

$$
(x-1,y+1).
$$

\textbf{Question:} After any number of operations, what can be concluded about $x+y$?
\end{frame}

\begin{frame}{Answer 4}
Initially,

$$
x+y=10.
$$

For the first operation,

$$
(x+1)+(y-1)=x+y=10.
$$

For the second operation,

$$
(x-1)+(y+1)=x+y=10.
$$

Thus, both possible operations preserve the sum.
\end{frame}

\begin{frame}{Answer 4}
Therefore, regardless of the number of operations,

$$
\boxed{x+y=10}.
$$

The quantity $x+y$ is an invariant of the program.

This allows us to reason about the program without examining every possible sequence of operations.
\end{frame}

% --------------------------------------------------
% Question 5
% --------------------------------------------------

\begin{frame}{Question 5}
A board contains tokens placed on grid positions.

Initially, the number of black tokens equals the number of white tokens.

An operation removes one black token and one white token simultaneously.

\textbf{Question:} Show that the difference between the number of black and white tokens is an invariant.
\end{frame}

\begin{frame}{Answer 5}
Let

$$
B=\text{number of black tokens}
$$

and

$$
W=\text{number of white tokens}.
$$

Initially,

$$
B=W.
$$

Therefore,

$$
B-W=0.
$$

\end{frame}

\begin{frame}{Answer 5}
After one operation,

$$
B\rightarrow B-1
$$

and

$$
W\rightarrow W-1.
$$

The new difference is

$$
(B-1)-(W-1)
=
B-W.
$$

Therefore, the difference does not change.
\end{frame}

\begin{frame}{Answer 5}
Hence,

$$
\boxed{B-W=0}
$$

remains true after every operation.

Therefore, the equality

$$
\boxed{B=W}
$$

is an invariant of the process.
\end{frame}
